\section{Divergence bridge}

This chapter expands the cycle bridge and the DWDB-DIV assembly.

\subsection{Cycle words}

A positive cycle gives a closed word
\begin{equation*}
  \wordOrbit(w,\Orbit_7(c))=\Orbit_7(c).
\end{equation*}
The word is split into rise steps ($t=1$ or $t=2$) and C3 steps
($t\ge3$).

\subsection{Construction of the period word}

Define the period word recursively by
\begin{equation*}
  \mathrm{cycleWord}(c,0)=[],
\end{equation*}
\begin{equation*}
  \mathrm{cycleWord}(c,p+1)=\mathrm{cycleWord}(c,p)
    ++ \left[\vtwo{5\,\Orbit_7(c+p)+1}\right].
\end{equation*}
By induction on $p$,
\begin{equation*}
  \wordOrbit(\mathrm{cycleWord}(c,p),\Orbit_7(c))
    =\Orbit_7(c+p),
\end{equation*}
and each appended entry is exact by definition.  Hence the compiled
lemmas \texttt{cycleWord\_step\_exact} and
\texttt{cycleWord\_take\_eq} hold.  If
$\Orbit_7(c)=\Orbit_7(c+p)$, the word is closed and
\texttt{cycleWord\_cycle\_equation} applies.

\subsection{Cycle word lemmas}

The following Lean lemmas extract the cycle word and its QB-8 split:
\begin{itemize}[leftmargin=*]
\item \texttt{orbit\_cycle\_imp\_min\_rise\_start};
\item \texttt{orbit\_reset\_head\_eq\_of\_rise\_start};
\item \texttt{orbit\_cycle\_imp\_min\_rise\_witness};
\item \texttt{cycleWord\_step\_exact};
\item \texttt{cycleWord\_take\_eq};
\item \texttt{cycleQb8Input\_exists\_rise\_index};
\item \texttt{cycleQb8Input\_exists\_c3\_index};
\item \texttt{cycleQb8Input\_exists\_block\_head\_mod\_five};
\item \texttt{rise\_block\_balance}.
\end{itemize}

\subsection{Extraction and split proof sketch}

Given a positive cycle, choose equal orbit values at depths $m<n$ and
put $c=m$, $p=n-m$, $m_0=\Orbit_7(c)$.  Define the word by
\begin{equation*}
  t_i=\vtwo{5\,\Orbit_7(c+i)+1},\qquad 0\le i<p.
\end{equation*}
The compiled lemmas \texttt{cycleWord\_step\_exact},
\texttt{cycleWord\_take\_eq}, and \texttt{cycleWord\_cycle\_equation}
give validity, closedness, and the exact cycle equation.

The QB-8 split is obtained by filtering the word into rise entries
$t<3$ and C3 entries $t\ge3$.  Closedness implies both filters are
nonempty: a cycle cannot consist only of rise steps, because their
total weight would be at most $2P$, contradicting the cycle equation
$2^Wm=5^Pm+A(w)$; and it cannot consist only of C3 steps, because each
such step satisfies $(5x+1)/2^t<x$ for $x\ge1$, so the states would
strictly decrease and could not close.  The precise statements are
\texttt{cycleQb8Input} and
\texttt{cycleQb8Input\_P\_ge\_two}.

\subsection{The QB-8 input record}

The Lean structure \texttt{CycleQb8Input m S P w rise c3} records
everything that follows from a positive cycle without range
assumptions:
\begin{itemize}[leftmargin=*]
\item \texttt{hvalid}: the word $w$ is valid from $m$;
\item \texttt{hclosed}: $\wordOrbit(w,m)=m$;
\item \texttt{hlength}, \texttt{hweight}: length $P$ and total
  weight $S$;
\item \texttt{hrise\_ok}, \texttt{hc3\_ok}: rise entries lie in
  $\{1,2\}$ and C3 entries are at least $3$;
\item \texttt{hlen\_split}, \texttt{hweight\_split}: the rise and C3
  lists decompose the word and its weight;
\item \texttt{hrise\_filter}, \texttt{hc3\_filter}: the split is the
  exact filter split;
\item \texttt{hc3\_pos}, \texttt{hrise\_pos}: both parts are
  nonempty;
\item \texttt{hexact}: every word step equals the orbit valuation;
\item \texttt{hm\_pos}, \texttt{hm\_odd}, \texttt{hm\_not\_five}:
  $m$ is positive, odd, and not divisible by $5$.
\end{itemize}

The structure is produced by
\texttt{orbit\_cycle\_imp\_cycle\_qb8\_input}.

\subsection{Failure window fields}

The structure \texttt{FailureWindow m S P w rise c3 j k0 t delta s}
adds to the QB-8 input:
\begin{itemize}[leftmargin=*]
\item \texttt{hj\_lt}: $j<P$;
\item \texttt{ht}: $t\in\{1,2\}$;
\item \texttt{hdelta}: $\delta=1$ when $t=1$, and
  $\delta\in\{1,3\}$ when $t=2$;
\item \texttt{hreset}: the exact reset equation with the prefix state
  $\wordOrbit(w{\upharpoonright}j,m)$;
\item \texttt{hreach}: the corrected reachability predicate
  $\ResetWindow(j,k_0,t,\delta,s)$;
\item \texttt{hs\_odd}, \texttt{hs\_not\_five}, \texttt{hk},
  \texttt{hs\_lt}: parity, coprimality, and size conditions;
\item \texttt{hfail\_t1}, \texttt{hfail\_t2}: the valuation lower
  bounds $2j+12$ and $2j+11$.
\end{itemize}

\subsection{Construction checklist}

The intended construction from a closed QB-8 word is conditional on
the unified core; the statement \texttt{failureWindowExistence} is not
claimed as proved in Lean.
\begin{enumerate}[leftmargin=*]
\item Extract the closed word and its rise/C3 split.
\item Choose the block head from the rise and C3 position lemmas.
\item Obtain the reset equation and the corrected reachability witness.
\item Apply the unified core to the block tail.
\item Read off the valuation lower bound.
\item Compare with the corrected decisive window bound.
\end{enumerate}

\subsection{Corrected window bound}

The corrected decisive window bound requires
$\ResetWindow(j,k_0,t,\delta,s)$ and asserts
\begin{equation*}
  \vtwo{5^{k_0+1}s+\delta5^j}\le2j+10
\qquad(t=2),
\end{equation*}
\begin{equation*}
  \vtwo{5^{k_0+1}s+5^j-2}\le2j+11
\qquad(t=1).
\end{equation*}

The old unconstrained bound is false; the counterexample is
\begin{equation*}
  j=36,\ k_0=0,\ t=2,\ \delta=3,\
  s=2^{83}-3\cdot5^{35},
\end{equation*}
with valuation $83>82$.  The corrected predicate
$\ResetWindow(j,k_0,t,\delta,s)$ fixes the exact block parameters and
carries full orbit reachability.

The valuation is 83 because the $5$-adic parts cancel exactly:
\begin{equation*}
  5s+3\cdot5^{36}
    =5(2^{83}-3\cdot5^{35})+3\cdot5^{36}
    =5\cdot2^{83}.
\end{equation*}
The witness satisfies the size, parity, and coprimality conditions;
it fails only because it carries no orbit reachability data.

\subsection{Corrected bound as a projection}

The corrected bound is not introduced as a separate analytic
hypothesis.  In the Lean repository it is recorded as a definition
with the unified core as an input:
\begin{itemize}[leftmargin=*]
\item \texttt{unifiedCoreClosed}: the unified-core valuation
  inequality under the no-$H_{\ge}$ premises, reachability, and
  $2\le H_s$;
\item \texttt{decisiveWindowValuationBoundCorrected\_of\_unified\_core}:
  \texttt{unifiedCoreClosed} $\to$
  \texttt{BlockAutomaton.decisiveWindowValuationBoundCorrected};
\item \texttt{failureWindowExistenceOfUnifiedCore}:
  \texttt{unifiedCoreClosed} $\to$
  \texttt{CycleBridge.failureWindowExistence}.
\end{itemize}

These are conditional interfaces.  None of them claims that the
unified core is already closed in Lean.

\subsection{Projection field mapping}

\begin{center}
\small
\begin{tabular}{@{}p{0.34\textwidth}p{0.30\textwidth}p{0.30\textwidth}@{}}
\toprule
\texttt{ResetWindowReachability} field & Role & Unified-core input \\
\midrule
$s5^{k_0}=r+1$ & previous terminal & block-tail terminal \\
$\GeneralOrbit(r)$ & terminal reachability & general-orbit input \\
$\FullOrbit(r_j)$ & reset-head reachability & block-head reachability \\
$\mathrm{ResetHeadEq}$ & exact reset equation & block equation \\
$s<5^{j-1-k_0}$, oddness, $5\nmid s$ & size and parity & candidate size \\
\bottomrule
\end{tabular}
\end{center}

\subsection{Corrected bound statement shape}

\begin{verbatim}
def t2WindowBoundCorrected (j k0 a t delta s : Nat) : Prop :=
  a = j - k0 - 1 ->
  t = 2 -> delta = 1 \/ delta = 3 ->
  S6Audit.ResetWindowReachability j k0 t delta s ->
  t2WindowValue j k0 delta s <= 2 * j + 10

def t1WindowBoundCorrected (j k0 a t s : Nat) : Prop :=
  a = j - k0 - 1 -> t = 1 ->
  S6Audit.ResetWindowReachability j k0 t 1 s ->
  t1WindowValue j k0 s <= 2 * j + 11

def decisiveWindowValuationBoundCorrected : Prop :=
  (forall j k0 a t delta s,
     t2WindowBoundCorrected j k0 a t delta s) /\
  (forall j k0 a t s,
     t1WindowBoundCorrected j k0 a t s)
\end{verbatim}

\subsection{No cycle implies unbounded: full proof}

Assume that the orbit of $x$ is bounded by $B$.  Then every value lies
in the finite set $\{0,\dots,B\}$, so the pigeonhole principle gives
indices $m<n$ with $\Orbit_x(m)=\Orbit_x(n)$.  This is a positive
cycle.  Contrapositively, the absence of a positive cycle implies that
the orbit is unbounded.  The statement is
\texttt{unbounded\_of\_no\_cycle}.

\subsection{Failure windows}

A failure window is a reachable reset block whose valuation lower bound
exceeds the corrected upper bound.  Its existence is
\texttt{failureWindowExistence}, an input to the assembly.

\subsection{Assembly}

\begin{equation*}
  \text{corrected window bound}\ \Rightarrow\ \neg\OrbitCycle(7)
  \ \Rightarrow\ \mathrm{IsUnboundedOrbit}(7).
\end{equation*}

The Lean interfaces are:
\begin{itemize}[leftmargin=*]
\item \texttt{CycleBridge.windowBoundToNoCycle};
\item \texttt{CycleBridge.failure\_window\_contradicts\_window\_corrected};
\item \texttt{DwdbDiv.dwdbDivFinalAssembly}.
\end{itemize}

All downstream statements that consume the corrected window bound or
the failure-window existence are conditional interfaces; they do not
claim those inputs are already proved.

\subsection{Failure window contradiction}

For $t=2$, the failure window gives
\begin{equation*}
  \vtwo{5^{k_0+1}s+\delta5^j}\ge2j+11,
\end{equation*}
while the corrected bound gives
\begin{equation*}
  \vtwo{5^{k_0+1}s+\delta5^j}\le2j+10.
\end{equation*}
For $t=1$, the lower bound is $2j+12$ and the upper bound is $2j+11$.
Both are contradictions.  The Lean theorem is
\texttt{failure\_window\_contradicts\_window\_corrected}.

\subsection{Contradiction lemma shapes}

The two elementary contradiction lemmas are:
\begin{verbatim}
theorem failure_window_contradicts_t2
    (j k0 delta s v : Nat)
    (hfail : 2 * j + 11 <= v)
    (hwin : v <= 2 * j + 10) : False

theorem failure_window_contradicts_t1
    (j k0 s v : Nat)
    (hfail : 2 * j + 12 <= v)
    (hwin : v <= 2 * j + 11) : False
\end{verbatim}

\subsection{DwdbDiv assembly}

The final assembly theorems are:
\begin{itemize}[leftmargin=*]
\item \texttt{DwdbDiv.dwdbDivFinalAssembly};
\item \texttt{DwdbDiv.five\_x\_plus\_one\_diverges\_at\_7\_of\_window\_bound};
\item \texttt{DwdbDiv.cycle\_of\_window\_bound\_contradiction};
\item \texttt{DwdbDiv.five\_x\_plus\_one\_diverges\_at\_7\_of\_failure\_window}.
\end{itemize}

These theorems are conditional on the corrected window bound and on
\texttt{failureWindowExistence}; neither input is claimed as already
proved.
